\chapter{Normalización}
\section{¿Qué es una norma?}
Indica cómo deber ser un producto o cómo debe funcionar un servicio para que sea seguro y responda a lo que el consumidor espera de él.

\section{Proceso Regulatorio: ¿qué y quién en España?}
El origen de la Asociación Española de Normalización, UNE, y de AENOR, se encuentra en la Asociación Española de Normalización y Certificación, creada en 1986.

En 2017 se ha procedido a un desdoblamiento de sus actividades por el cual UNE, asociación sin fines lucrativos, desarrolla la actividad de Normalización y Cooperación. Por su parte, AENOR, entidad mercantil, trabaja en los ámbitos de la evaluación de la conformidad y actividades asociadas, como la formación o la venta de publicaciones.

La Asociación Española de Normalización, UNE, es el organismo legalmente responsable del desarrollo y difusión de las normas técnicas en España. Las normas indican cómo deber ser un producto o cómo debe funcionar un servicio para que sea seguro y responda  alo que el consumidor espera de él. UNE pone a disposición de todos uno de los catálogos más completos, con más de 31.500 documentos normativos que contienen soluciones eficaces.

UNE aporta su experiencia y conocimiento en materia de normas y de productos y servicios relacionados a organizaciones de todo el mundo, desarrollando una gran actividad de cooperación internacional.

En evaluación de la conformidad, el trbajo serio y riguroso que caracteriza el trabajo de la Entidad desde su creación ha posibilitado que los certificados de AENOR sean los más valorados, no sólo en España sino también en el ámbito internacional, habiendo emitido certificados en másd e 60 países. AENOR se sitúa entre las 10 certificadoras más importantes del mundo.

La certificación es la acción llevada a cabo por una entidad independiente de las partes interesadas mediante la que se manifiesta que una organización, producto, proceso o servicio, cumple los requisitos definidos en unas normas o especificaciones técnicas.

El Parlamento Europeo y el Consejo Europeo definieron y aprobaron un Reglamento europeo sobre normalización (Reglamento (UE nº 1025/2012)): publicado en el Boletín Oficial de la U.E. el 14 de noviembre de 2012 y siendo obligatorio para todos los Estados miembros a partir del 1 de enero de 2013.

Esta regulación define el marco legal de las Entidades Europeas de Estandarización:
\begin{itemize}
    \item CEN, el Comité Europeo de Normalización, es una asociación que reúne a los Organismos Nacionales de Normalización de 33 países europeos.
    \item CENELEC es el Comité Europeo de Normalización Electrotécnica y es responsable de la normalización en el campo de la ingnienería electrotécnica. CENELEC prepara normas voluntarias que ayudan a facilitar el comercio entre países, crear nuesvos mercados, reducir los costos de cumplimiento y apoyar el desarrollo de un mercado úncio europeo.
    \item ETSI, el Instituto Europeo de Normas de Telecomunicaciones, produce normas aplicables a nivel mundial para las Tecnologías de la Información y las Comunicaciones (TIC), incluidas las tecnologías fija, móvil, de radio, convergente, de radiodifusión e Internet. Las normas facilitan el trabajo de las tecnologías de las que dependen los negocios y la sociedad. Por ejemplo, los estándares para GSM, DECT, tarjetas inteligentes y firmas electrónicas han ayudado a revolucionar la vida moderna en todo el mundo.
\end{itemize}

A partir de esta estrategia y la relación de las tres organizaciones, los estándares se producen bajo un contexto legal que incluye directivas o regulaciones específicas.